\documentclass[12pt]{article}
\usepackage[margin=1in]{geometry}
\usepackage{amsmath}
\usepackage{natbib}
\usepackage{booktabs}
\usepackage{hyperref}

\title{Sequestered Debts: Virginia, the British Merchant Debts,\\
and the Reputation of the New Republic}
\author{George J. Hall and Thomas J. Sargent\thanks{Background memorandum
prepared with Claude for the authors' fiscal-history project. Claims marked
\texttt{[PRIMARY]} in the accompanying \texttt{virginia\_british\_debts.bib}
were checked against primary texts---the statute as quoted in \emph{Ware v.
Hylton}, the treaty and convention texts on the Avalon Project, and the U.S.
Reports. Several figures could not be independently verified; they are collected
in the caveats of Section~\ref{sec:caveats}.}}
\date{\today}

\begin{document}

\maketitle

\begin{abstract}
During the War of Independence Virginia sequestered the debts its citizens owed
to British merchants and let those debtors extinguish sterling obligations by
paying depreciated paper into the state loan office at par---a windfall for the
most heavily indebted planters in America. The Treaty of Paris (1783) promised
British creditors ``no lawful impediment'' to recovering the full sterling value
of such debts, but the states obstructed it, and the debt question became
entangled in a decade-long mutual-breach standoff over the British-held frontier
posts and the enslaved people carried away at the war's end. The Supreme Court
resolved the legal question in \emph{Ware v. Hylton} (1796), holding that the
treaty, as supreme federal law, annulled the Virginia statute---the Court's first
assertion of treaty supremacy over state law, and John Marshall's only (losing)
appearance as an advocate. Practical recovery nonetheless required a political
settlement: the Jay Treaty's debt commission (1797--1799) deadlocked, and the
Convention of 1802 discharged the matter by a federal lump-sum payment of
\pounds600{,}000. Honoring these debts, alongside Hamilton's funding and the full
servicing of the Dutch and French loans, helped establish the young republic's
reputation as a sovereign that paid---reassuring European creditors and lowering
its cost of capital. Virginia's device was the archetype but not unique: North
Carolina, Maryland, and Georgia adopted variants, concentrated in the tobacco
Chesapeake and lower South where the British mercantile debt was largest.
\end{abstract}

\section{Introduction}

This memorandum assembles, with sources, an episode that bears directly on the
reputational themes of our larger project: how the United States came to honor---
after resisting---the private debts its citizens owed to British merchants, and
what that resolution signaled to the country's other creditors. The story runs
from a Virginia wartime statute of 1777, through the Treaty of Paris and a
constitutional crisis over treaty compliance, to a landmark Supreme Court
decision and, finally, a cash settlement in 1802. It is at once a debtor-relief
story, a constitutional story about the supremacy of federal treaties, and a
public-credit story about the value of a reputation for paying.

\section{Virginia's Sequestration of British Debts, 1777--1780}
\label{sec:virginia}

At its October 1777 session the Virginia General Assembly passed ``An act for
sequestering British property, enabling those indebted to British subjects to pay
off such debts'' \citep[9:377--380]{hening1777sequestration}. The bill was drafted
by Thomas Jefferson; because the session sat into the new year, the literature
dates it variously to October 1777 and January 1778---the same act.

The mechanism deserves precise statement, because it differs slightly from
``buying state bonds at a discount.'' Under the act's third section---quoted
verbatim in \emph{Ware v. Hylton}---any Virginian owing money to a British subject
could pay the sum into the state \emph{loan office}, take out a loan-office
certificate in the creditor's name, and deliver it to the Governor and Council,
whose receipt ``shall discharge him from so much of the debt''
\citep{warevhylton1796, hening1777sequestration}. The instrument the debtor
received was thus a Virginia loan-office certificate---a state obligation, or
``state bond,'' issued in the British creditor's name---and the Commonwealth
interposed itself between debtor and creditor. Notably, the statute said nothing
about whether Virginia was itself obliged to pay the British creditor.

The discount lay in the currency. Debtors paid in depreciating Virginia and
Continental paper, credited at face (par) value against debts contracted in
sterling ``good British money.'' Because that paper was collapsing---and because
the paper and the certificates themselves traded far below face---a planter could
retire a hard-money sterling obligation for a small fraction of its real value,
the economic equivalent of buying the medium of payment at a deep discount and
tendering it at par. The case record supplies a vivid illustration: the Hyltons'
bond to the Bristol firm Farell \& Jones, dated 7~July 1774, was for
\pounds2{,}976~11s.~6d.\ sterling; on 26~April 1780 the debtors paid
\pounds933~14s.\ in Virginia currency (\$3{,}111\ensuremath{\tfrac{1}{9}}) into the loan
office---a nominally smaller sum that, in specie terms at that stage of the
depreciation, was worth only a sliver of the sterling debt
\citep{warevhylton1796}. The pay-in privilege was repealed in May 1780 as the
currency disintegrated.

The stakes were largest in Virginia because Chesapeake tobacco planters were the
most indebted colonists in America. Of roughly \pounds5 million owed by the
thirteen colonies to British merchants on the eve of the Revolution, Virginia's
share was about \pounds2.3 million---Jefferson reckoned Virginians owed ``over
\pounds2 million'' in 1774, perhaps half of all American private
indebtedness, most of it retail store credit to Glasgow and English firms
\citep{adamson2003british, evans1962planter, evans1971private}. That the debt was
so widely distributed among planters made it politically explosive
\citep{holton1999forced, sloan1995principle}.

\section{The Treaty of Peace and the Debt Question, 1783}
\label{sec:treaty}

The Definitive Treaty of Peace (Paris, 3~September 1783) confronted these state
laws head-on. \emph{Article 4} provided \citep{treatyparis1783}:
\begin{quote}
It is agreed that creditors on either side shall meet with no lawful impediment to
the recovery of the full value in sterling money of all bona fide debts heretofore
contracted.
\end{quote}
The single sentence repudiated the state debtor-relief statutes: it promised
British creditors the \emph{full sterling value}, notwithstanding any law that had
let debtors discharge such debts in depreciated paper. Two neighboring articles
rounded it out. \emph{Article 5} had Congress only ``earnestly recommend'' to the
states the restitution of confiscated Loyalist estates---a recommendation Congress
was powerless to compel---and \emph{Article 6} barred future confiscations and
prosecutions for wartime conduct \citep{treatyparis1783}.

The promise then dissolved into a mutual-breach standoff that lasted a decade. The
states---Virginia foremost---kept their obstructing laws and effectively closed
their courts to British creditors; Britain in turn refused to evacuate the
Northwest frontier posts (Detroit, Niagara, and others) required by Article~7,
and at the 1783 evacuation of New York, Sir Guy Carleton allowed roughly $3{,}000$
formerly enslaved people (recorded in the ``Book of Negroes'') to be carried away,
which Southern planters demanded compensation for as its own violation of
Article~7 \citep{ritcheson1969aftermath}. Each side blamed the other for the first
breach. Secretary for Foreign Affairs John Jay, reporting to Congress on
13~October 1786, concluded that the American \emph{states} had breached
first through their debt and confiscation laws \citep{jayreport1786}; on his
recommendation Congress resolved (21~March 1787) that treaties are ``the supreme
law of the land'' binding on the states and that all repugnant state acts ``ought
to be forthwith repealed'' \citep{congress1787resolve}. The episode is a
documented motive for the Constitution's Supremacy Clause and treaty power
\citep{edling2003revolution}. In Virginia, Patrick Henry led the resistance, and
even after 1787 the Assembly agreed only to open its courts \emph{conditionally}
---on British evacuation of the posts and compensation for the enslaved people
carried away---so recovery remained blocked
\citep{madison1784resolutions, encvirginia_treaty}.

\section{The Supreme Court and Treaty Supremacy}
\label{sec:court}

\subsection{Ware v. Hylton (1796)}

The legal question was settled in \emph{Ware v. Hylton}, 3~U.S.\ (3~Dall.)\ 199
(1796)---``the British Debt Case.'' Ware, administrator of the estate of William
Jones (surviving partner of Farell \& Jones of Bristol), sued the Virginia debtors
Daniel Hylton \& Co.\ and Francis Eppes on the 1774 bond; the debtors pleaded the
\pounds933 they had paid into the loan office as a discharge. The Court held that
Article~4 of the 1783 treaty, as the supreme law of the land under Article~VI,
overrode and nullified the Virginia sequestration statute, so that the loan-office
payment did \emph{not} discharge the debt and the creditor could recover in full
with interest \citep{warevhylton1796}.

The Justices delivered opinions seriatim. Samuel Chase wrote the leading one:
conceding that Virginia had sovereign power in 1777 to sequester enemy debts, he
held that the treaty annulled the law's effect and revived the creditor's right,
state law being ``prostrate'' before a federal treaty. Paterson, Wilson, and
Cushing concurred in the result; Iredell alone reasoned for the debtors,
consistent with the ruling he had given below on circuit (his position is best
described substantively rather than as a formal dissent). Two features have lasted:
John Marshall argued for the debtors---his only appearance as an advocate before
the Supreme Court, and he lost---five years before he became Chief Justice; and the
decision was the Court's first assertion of treaty supremacy over conflicting
state law, an exercise of judicial review that preceded \emph{Marbury v. Madison}
by seven years \citep{warevhylton1796, hobson1996great, currie1985constitution,
goebel1971history}.

\subsection{The related cases}

\emph{Ware} had a circuit-court precursor, \emph{Jones v. Walker} (C.C.D.\ Va.\
1791), in which Patrick Henry argued for the debtor before Justice Iredell but
which was left undecided \citep{jonesvwalker1791, marshallpapers1987}. And
\emph{Georgia v. Brailsford} (1792--1794) tested Georgia's sequestration of a
British debt; its 1794 phase is famous as the only jury trial ever held by the
Supreme Court, and the Court held that sequestration merely \emph{suspended}
recovery during the war rather than extinguishing it, so that peace revived the
creditor's right \citep{georgiavbrailsford1794}. (A biographical thread worth
noting: Farell \& Jones were also the creditors of John Wayles, Jefferson's
father-in-law, so Jefferson himself was among the Virginians entangled in these
debts.)

\section{The Political Settlement, 1794--1802}
\label{sec:settlement}

\emph{Ware} established the creditors' legal right, but recovery through hostile
local juries remained so uncertain that the matter was resolved politically.
Article~6 of the Jay Treaty (19~November 1794) created a five-member commission
(two British, two American, and a fifth---John Guillemard, a British subject---
chosen by lot) under which the United States promised ``full and complete
compensation'' for pre-war debts that lawful impediments had made unrecoverable,
excluding losses due to debtor insolvency \citep{jaytreaty1794}. The commission
met in Philadelphia from May 1797 but broke down in July 1799, the American
commissioners withdrawing over an effectively 3--2 British-aligned majority on
questions of jurisdiction, wartime interest, and debtor solvency
\citep{bemis1923jay, ritcheson1969aftermath}.

The deadlock was cut by the Convention of 8~January 1802 (Rufus King and Lord
Hawkesbury), under which the United States agreed to pay Britain a lump sum of
\pounds600{,}000 sterling (about \$2.66 million), in three annual installments of
\pounds200{,}000, expressly ``cancelling and annulling'' Jay Treaty Article~6
\citep{convention1802}. Britain then distributed the money through its own
commission, which sat until 1811 and approved only a fraction of the claims
presented.

\section{Reputation and Relations with Other Creditors}
\label{sec:reputation}

The reputational payoff is the point that connects this episode to public credit.
By enforcing the treaty in \emph{Ware} and then paying the \pounds600{,}000 rather
than repudiating, the federal government signaled a sovereign that would honor both
its debts and its treaty obligations. Combined with Hamilton's funding and
assumption of 1790 and the full servicing and refinancing of the Dutch and French
loans, it helped give the young United States, within a few years, essentially the
highest credit rating in the world \citep{sylla2011financial, edling2014hercules}.
That reputation reassured the Amsterdam bankers holding U.S.\ paper---many of whom
had bought at deep discount and were now repaid in full---and the French; it
lowered borrowing costs and underwrote cheap access to European capital markets,
through to the London-and-Amsterdam financing of the Louisiana Purchase in 1803
\citep{statedept_loans, sylla2011financial}. Honoring the British merchant debts
was, in effect, the outward-facing counterpart to Hamilton's domestic funding:
both were investments in a reputation for paying.

\section{Were Other States Doing the Same Thing?}
\label{sec:other-states}

Virginia was the archetype but not unique. \emph{Ware v. Hylton} itself records
that ``in several other states, laws had been passed authorising a discharge of
British debts in paper money, or by a tender of property,'' and the standard
reference notes that ``all of the newly declared states enacted laws that impaired
the position of creditors'' \citep{warevhylton1796, adamson2003british}. It helps
to distinguish three mechanisms, summarized in Table~\ref{tab:states}: (a) the
debtor pays the British debt into the state and is discharged in depreciated paper
(the Virginia model); (b) outright confiscation of the debt as enemy property; and
(c) mere suspension of recovery or closure of the courts.

\begin{table}[htbp]
\centering
\caption{State treatment of debts owed to British merchants, 1776--1783}
\label{tab:states}
\small
\begin{tabular}{p{2.4cm} p{2.2cm} p{8.3cm}}
\toprule
State & Mechanism & Description \\
\midrule
Virginia & (a) pay-in & 1777 act: debtor pays depreciated paper into the loan office at par, takes a certificate in the creditor's name, and is discharged; repealed 1780 \\
North Carolina & (a)/(b) hybrid & 1777 act: debtor owing sterling to a confiscated person pays \pounds175 N.C.\ currency per \pounds100 sterling into the state treasury \\
Maryland & (b) estates; (a) in practice & 1780 act confiscated Loyalist \emph{estates} (British debts exempted by the Senate), yet an official 1781 ``black list'' records debtors who paid depreciated currency into the treasury on British debts \\
Georgia & (c) suspension & Sequestered British debts, barring recovery by suit during the war; peace and the treaty revived the right (\emph{Georgia v. Brailsford}) \\
South Carolina & (b) estates & Confiscation and amercement of Loyalist estates (1782); a debtor pay-in is not documented \\
New York, Pennsylvania, etc. & (b) estates & Widespread confiscation of Loyalist \emph{estates} (e.g.\ New York's Forfeiture and Trespass Acts), but not the British-merchant-debt pay-in \\
\bottomrule
\end{tabular}
\end{table}

The pattern is regional and follows the money: the discharge-by-pay-in device
clustered in the tobacco Chesapeake and lower South, precisely where the British
mercantile debt was concentrated---Virginia alone owed roughly \pounds2.3 million
of the \pounds5 million total \citep{adamson2003british, kellock1974london}. It
was this thicket of state laws---several allowing paper discharge, more merely
obstructing recovery---that Article~4 of the Treaty of Paris, then \emph{Ware v.
Hylton}, then the Jay Treaty commission, and finally the \pounds600{,}000
settlement of 1802 were successively deployed to override.

\section{Caveats and Verification Notes}
\label{sec:caveats}

Several specifics in this memorandum rest on secondary or reference sources and
should be confirmed against primary materials before they are woven into the book:

\begin{itemize}
  \item \textbf{The aggregate paid into the Virginia loan office.} A total of
    roughly \pounds273{,}000 is sometimes cited, but we could not verify it from an
    accessible source; it is most likely drawn from \citet[349--360]{evans1971private}
    and should be checked there. The single-case figures (the \pounds2{,}976 bond
    and the \pounds933 payment) are from the \emph{Ware} record and are secure.
  \item \textbf{The Maryland pay-in.} That Maryland let debtors discharge British
    debts in depreciated paper is supported by \emph{Ware}'s ``several other
    states'' language and by the 1781 ``black list'' controversy, but the exact
    statute, its terms, and whether the pay-in was mandatory are not confirmed;
    note also that the Maryland Senate exempted British debts from the 1780
    confiscation act, so Maryland's treatment was more tangled than Virginia's.
    Verify in the \emph{Archives of Maryland}.
  \item \textbf{The plaintiff's full name.} The reports caption the plaintiff simply
    ``Ware, Administrator of Jones.'' The administrator is commonly given as ``John
    Tyndale Ware,'' but we did not confirm the full name from a primary source.
  \item \textbf{Spelling of the creditor firm.} The Bristol house appears as both
    ``Farell \& Jones'' and ``Farrell \& Jones''; Founders Online uses ``Farell.''
  \item \textbf{Iredell's posture in \emph{Ware}.} Iredell had favored the debtors
    on circuit and his Supreme Court opinion reprised that view; whether to label it
    a ``dissent'' varies among sources, so we describe it substantively.
  \item \textbf{Number of obstructing states.} Sources say ``all'' the states
    impaired creditors in some way and ``several'' allowed the paper discharge
    specifically; no reliable source gives a precise count.
  \item \textbf{Monograph page numbers.} Specific pages in Bemis, Ritcheson, Sloan,
    and Edling for the \pounds600{,}000 settlement and the reputational argument were
    not pinned to the physical volumes and should be checked before quotation.
\end{itemize}

\bibliographystyle{aer}
\bibliography{virginia_british_debts}

\end{document}
